\documentclass[11pt,aspectratio=169]{beamer}

% ================= Theme: Metropolis (clean, modern, flat) =================
\usetheme{metropolis}

% ---------- Custom "water / irrigation" colour palette ----------
% Deep teal primary, warm amber accent, soft off-white background
\definecolor{deepteal}{HTML}{0B4F6C}
\definecolor{amber}{HTML}{E8A33D}
\definecolor{leaf}{HTML}{3B8C6E}
\definecolor{offwhite}{HTML}{FAFAF7}
\definecolor{slategray}{HTML}{4A5A5F}

\setbeamercolor{palette primary}{bg=deepteal, fg=white}
\setbeamercolor{frametitle}{bg=deepteal, fg=white}
\setbeamercolor{progress bar}{fg=amber, bg=slategray!20}
\setbeamercolor{alerted text}{fg=amber}
\setbeamercolor{normal text}{bg=offwhite}
\setbeamercolor{background canvas}{bg=offwhite}
\setbeamercolor{block title}{bg=leaf, fg=white}
\setbeamercolor{block body}{bg=leaf!10, fg=black}
\setbeamertemplate{blocks}[rounded]

\setbeamertemplate{navigation symbols}{}
\metroset{sectionpage=progressbar, subsectionpage=progressbar, numbering=fraction, progressbar=frametitle, block=fill}

\usepackage[utf8]{inputenc}
\usepackage{graphicx}
\usepackage{booktabs}
\usepackage{array}
\usepackage{multirow}

% shrink default vertical gaps a touch so bullet+figure slides fit safely
\setlength{\parskip}{2pt}

\title{Predicting Irrigation Water Demand}
\subtitle{with Machine Learning}
\author{Akhil Nistala}
\institute{Civil Engineering (Water Resources Engineering) Term Project \\ IIT Hyderabad}
\date{\today}

% Helper: a figure that is guaranteed to fit under a frame title + caption,
% no matter its native aspect ratio.
\newcommand{\fitfig}[2][0.62]{%
  \includegraphics[width=\textwidth,height=#1\textheight,keepaspectratio]{#2}%
}

\begin{document}

% ---------------- Title ----------------
{
\setbeamertemplate{footline}{}
\begin{frame}
  \titlepage
\end{frame}
}

% ---------------- Outline ----------------
\begin{frame}{Overview}
  \tableofcontents
\end{frame}

\section{The Problem}

\begin{frame}{Why does this matter?}
  \begin{itemize}
    \item Every day, a farmer must decide: \alert{how much water to give the crop?}
    \item Too little $\rightarrow$ crop stress, lower yield.
    \item Too much $\rightarrow$ wasted water, waterlogging, nutrient loss.
    \item Engineers already know how to calculate this correctly --- but it needs
      five different weather readings and a fairly involved calculation, every
      single day, for every field.
  \end{itemize}
  \begin{alertblock}{The question this project asks}
    Can a machine-learning model \textbf{learn} to predict the same answer directly
    from simple weather and crop data --- fast, without redoing the full
    calculation each time?
  \end{alertblock}
\end{frame}

\section{The Idea in Plain Terms}

\begin{frame}{Think of the field as a bucket}
  \begin{columns}[T]
    \begin{column}{0.55\textwidth}
      \begin{itemize}
        \item Water \textbf{comes in} from rain and irrigation.
        \item Water \textbf{leaves} through evaporation from the soil and
          transpiration by the plant (together, ``evapotranspiration'').
        \item If more leaves than comes in, the bucket runs low --- and
          irrigation must top it up.
      \end{itemize}
      \vspace{0.5em}
      That's the entire physical idea behind this project.
    \end{column}
    \begin{column}{0.42\textwidth}
      \centering
      \small
      \textcolor{deepteal}{\textbf{Rain + Irrigation}} \\[0.3em]
      $\big\Downarrow$ \\[0.3em]
      \textcolor{leaf}{\textbf{Soil--Crop Bucket}} \\[0.3em]
      $\big\Uparrow$ \\[0.3em]
      \textcolor{amber!80!black}{\textbf{Evapotranspiration}}
    \end{column}
  \end{columns}
\end{frame}

\begin{frame}{The standard engineering method: FAO-56}
  \begin{itemize}
    \item FAO-56 is the internationally accepted method for estimating how much
      water a crop needs, published by the UN Food and Agriculture Organization.
    \item It works in two simple steps:
  \end{itemize}
  \begin{center}
    \small
    \textbf{Weather} $\;\rightarrow\;$ \textbf{Atmospheric thirst (ET0)} \\[0.2em]
    $\times$ \\[0.2em]
    \textbf{Crop coefficient (Kc)} \; -- depends on crop type \& growth stage \\[0.2em]
    $=$ \\[0.2em]
    \textbf{Actual crop water need}
  \end{center}
  Then subtract whatever rainfall already supplied, and divide by irrigation
  system efficiency, to get the final \alert{irrigation demand}.
\end{frame}

\section{The Data}

\begin{frame}{Where did the data come from?}
  \begin{itemize}
    \item No real farm dataset with all the needed variables was available.
    \item So a \textbf{synthetic dataset} was generated --- not random
      numbers, but the exact FAO-56 physics above, run over realistic Indian
      weather patterns.
    \item Weather behaves like real weather: hot spells and wet spells last
      several days, not one isolated day, and rainfall follows a realistic
      ``many small showers, occasional big storm'' pattern.
    \item \textbf{2,000 days} (about 5.5 years) of daily records, for a maize
      crop, at a representative central-Indian location.
    \item Later, this weather is replaced with \textbf{real NASA POWER data}
      for the same place and dates.
  \end{itemize}
\end{frame}

\begin{frame}{What goes into the model}
  \begin{columns}[T]
    \begin{column}{0.5\textwidth}
      \textbf{Inputs (what we know):}
      \begin{itemize}
        \item Rainfall
        \item Temperature
        \item Humidity
        \item Wind speed
        \item Solar radiation
        \item Crop coefficient
        \item Growth stage
      \end{itemize}
    \end{column}
    \begin{column}{0.5\textwidth}
      \textbf{Output (what we predict):}
      \begin{itemize}
        \item Today's irrigation water demand (mm/day)
      \end{itemize}
      \vspace{1em}
      \begin{block}{}
        \small The model never sees the intermediate physics calculations ---
        only the raw inputs a farmer would actually have.
      \end{block}
    \end{column}
  \end{columns}
\end{frame}

\begin{frame}{What the data looks like}
  \centering
  \fitfig[0.68]{figures/eda_grid.png}
  \vspace{0.3em}

  \small Hotter, sunnier days $\rightarrow$ more water needed. Rainy days $\rightarrow$ demand drops to near zero.
\end{frame}

\section{Building the Models}

\begin{frame}{Splitting the data fairly}
  \begin{itemize}
    \item First \textbf{80\% of days} (chronologically) used to train the
      models.
    \item Last \textbf{20\% of days} held back to test them.
    \item This mimics real life: a model only ever has the past to predict
      the future --- it never gets to ``peek'' at days close to the ones it's
      tested on.
  \end{itemize}
\end{frame}

\begin{frame}{Three models compared}
  \begin{itemize}
    \item \textbf{Linear Regression} --- the simplest approach: assumes a
      straight-line relationship between inputs and demand.
    \item \textbf{Decision Tree} --- learns a series of yes/no questions
      (``Is it hot? Did it rain?'') to arrive at a prediction.
    \item \textbf{Random Forest} --- builds hundreds of decision trees and
      averages their answers, cancelling out individual mistakes.
  \end{itemize}
  \vspace{0.5em}
  These are then compared against three \textbf{deep-learning} models, first
  on synthetic weather and then on real NASA weather.
\end{frame}

\section{Results}

\begin{frame}{How good were the predictions?}
  \begin{table}
    \centering
    \small
    \begin{tabular}{lccc}
      \toprule
      \textbf{Model} & \textbf{MAE} & \textbf{RMSE} & \textbf{R\textsuperscript{2}} \\
       & \footnotesize (mm/day) & \footnotesize (mm/day) & \footnotesize (variance explained) \\
      \midrule
      Linear Regression & 0.83 & 1.11 & 91\% \\
      Decision Tree & 0.64 & 1.03 & 92\% \\
      \textbf{Random Forest} & \textbf{0.40} & \textbf{0.70} & \textbf{96\%} \\
      \bottomrule
    \end{tabular}
  \end{table}
  \footnotesize
  \textbf{MAE} = typical size of the daily prediction error, in mm/day. \quad
  \textbf{RMSE} = similar, but penalises the occasional big miss more heavily.
  \normalsize
  \begin{itemize}
    \item Random Forest was clearly the best --- it explains about 96\% of the
      day-to-day variation in irrigation demand.
    \item Linear Regression's straight-line assumption cannot capture the
      real, more complex weather-to-demand relationships.
  \end{itemize}
\end{frame}

\begin{frame}{A visual check: predicted vs.\ actual}
  \centering
  \fitfig[0.62]{figures/09b_lr_vs_rf_actual_vs_predicted.png}

  \small Points closer to the dashed line = better predictions. Linear
  Regression (left) predicts \textbf{negative} water demand on 13 of the 400
  test days --- physically impossible. Random Forest (right) never does.
\end{frame}

\begin{frame}{Model comparison at a glance}
  \centering
  \fitfig[0.6]{figures/09_classical_models.png}
\end{frame}

\section{What Matters Most}

\begin{frame}{What did the model rely on most?}
  \begin{itemize}
    \item A single ``importance'' score can be misleading, so \textbf{three
      independent measures} were checked against each other:
    \begin{enumerate}
      \item \textbf{Correlation} with the target --- simplest baseline.
      \item \textbf{Random Forest importance} --- how often/effectively a
        feature was used to build the trees.
      \item \textbf{Permutation importance} --- scramble one feature at a
        time on unseen data and see how much accuracy drops.
    \end{enumerate}
    \item The two Random Forest measures \alert{agree on the top drivers}.
      Simple correlation is misleading on its own: growth stage correlates
      with demand but adds nothing once the crop coefficient is known.
  \end{itemize}
\end{frame}

\begin{frame}{Feature importance --- what drives the prediction}
  \centering
  \fitfig[0.56]{figures/12_feature_importance_no_lag.png}

  \small
  \begin{itemize}
    \item \textbf{Crop coefficient} and \textbf{solar radiation} lead,
      followed by rainfall and temperature: the same drivers as the FAO-56
      method itself.
    \item Growth stage adds nothing extra, because the crop coefficient
      already encodes it.
  \end{itemize}
\end{frame}

\begin{frame}{Does the seasonal pattern make sense?}
  \centering
  \fitfig[0.62]{figures/08_monthly_avg_demand.png}

  \small Peaks in the hot, dry months before the monsoon (March--May); drops
  sharply during the monsoon, when rain does the work instead --- exactly
  what farmers and engineers already expect.
\end{frame}

\section{Comparison with Deep Learning}

\begin{frame}{Three deep-learning models}
  \small
  \begin{itemize}
    \item \textbf{MLP} (multi-layer perceptron): a small neural network
      on today's 7 inputs (2 hidden layers of 64, $\approx$4.7k weights).
    \item \textbf{1D-CNN}: scans the \textbf{last 7 days} of the same 7 inputs
      for short weather patterns ($\approx$18k weights).
    \item \textbf{LSTM}: a recurrent network that reads the last 7 days
      in order and keeps a memory of them ($\approx$12.5k weights).
  \end{itemize}
  \begin{block}{Kept fair}
    \footnotesize
    \begin{itemize}
      \item Same 7 inputs and the same 400 test days. The sequence models see
        past \emph{weather}, never past demand.
      \item Early stopping uses the last 10\% of the training period; the test
        days are never touched.
      \item Each network is trained with 3 random seeds and the average is
        reported.
      \item Control: a Random Forest given the same 7-day window.
    \end{itemize}
  \end{block}
\end{frame}

\begin{frame}{Classical ML vs.\ deep learning (7 inputs)}
  \begin{table}
    \centering
    \small
    \begin{tabular}{llccc}
      \toprule
      & \textbf{Model} & \textbf{MAE} & \textbf{RMSE} & \textbf{R\textsuperscript{2}} \\
      \midrule
      \multirow{4}{*}{\rotatebox{90}{\footnotesize Classical}}
      & Linear Regression & 0.83 & 1.11 & 0.908 \\
      & Decision Tree & 0.64 & 1.03 & 0.921 \\
      & Random Forest & 0.40 & 0.70 & 0.963 \\
      & Random Forest + 7-day window & 0.55 & 0.90 & 0.940 \\
      \midrule
      \multirow{3}{*}{\rotatebox{90}{\footnotesize Deep}}
      & MLP (today only) & 0.28 {\scriptsize$\pm$0.03} & 0.37 & 0.990 \\
      & 1D-CNN (7-day) & 0.35 {\scriptsize$\pm$0.02} & 0.49 & 0.982 \\
      & \textbf{LSTM (7-day)} & \textbf{0.26} {\scriptsize$\pm$0.01} & \textbf{0.36} & \textbf{0.990} \\
      \bottomrule
    \end{tabular}
  \end{table}
  \footnotesize MAE, RMSE in mm/day on the 400 test days; deep models are the
  mean of 3 seeds ($\pm$ 1 standard deviation).
  \small
  \begin{itemize}
    \item The best network (LSTM) cuts the typical error by \alert{35\%}
      versus Random Forest (0.40 $\rightarrow$ 0.26 mm/day), and the large
      misses (RMSE) by about half.
  \end{itemize}
\end{frame}

\begin{frame}{All models at a glance}
  \centering
  \fitfig[0.72]{figures/10_ml_vs_dl_no_lag.png}
\end{frame}

\begin{frame}{Predicted vs.\ actual: best tree model vs.\ best network}
  \centering
  \fitfig[0.6]{figures/11_actual_vs_predicted_no_lag.png}

  \small Random Forest predicts in ``steps'' and has a few large misses on
  rainy days. The LSTM follows the 1:1 line closely, though it occasionally
  predicts slightly below zero on dry-spell days (fixed by clipping at 0).
\end{frame}

\begin{frame}{Is the LSTM overfitting?}
  \begin{table}
    \centering
    \small
    \begin{tabular}{cccc}
      \toprule
      \textbf{Seed} & \textbf{Training days} & \textbf{Validation days} & \textbf{Test days} \\
       & \footnotesize MAE / R\textsuperscript{2} & \footnotesize MAE / R\textsuperscript{2} & \footnotesize MAE / R\textsuperscript{2} \\
      \midrule
      0 & 0.18 / 0.995 & 0.20 / 0.995 & 0.25 / 0.991 \\
      1 & 0.19 / 0.994 & 0.19 / 0.996 & 0.27 / 0.990 \\
      2 & 0.19 / 0.995 & 0.20 / 0.996 & 0.27 / 0.990 \\
      \bottomrule
    \end{tabular}
  \end{table}
  \vspace{-0.6em}
  \footnotesize
  \begin{itemize}
    \item \textbf{Small gap:} test error is only $\approx$0.07 mm/day higher than
      training error; R\textsuperscript{2} drops just 0.995 $\rightarrow$ 0.990.
    \item \textbf{Not memorising noise:} the data contains random noise worth
      $\approx$0.14 mm/day of error that no honest model can remove. Training
      error (0.19) stays above this floor.
    \item \textbf{Early stopping worked:} training stopped at epochs 98--118,
      when validation error stopped improving; all 3 seeds agree.
    \item The small remaining gap is expected: test days are the
      \emph{latest} 400 days, later than anything seen in training.
  \end{itemize}
  \alert{Verdict: no meaningful overfitting.}
\end{frame}

\begin{frame}{What the deep-learning comparison tells us}
  \small
  \begin{itemize}
    \item \textbf{Why networks win here:} irrigation demand is a smooth,
      multiplicative formula (crop coefficient $\times$ atmospheric thirst,
      minus rain). Neural networks fit smooth curves; trees can only
      approximate them in steps.
    \item \textbf{History adds little:} the MLP sees only today, yet is almost
      as good as the 7-day LSTM (0.28 vs 0.26 mm/day). Today's weather
      already carries most of the information.
    \item Giving the Random Forest the 7-day window actually made it
      \alert{worse} (0.40 $\rightarrow$ 0.55): 49 inputs instead of 7 dilute
      the useful ones.
    \item \textbf{Cost:} the networks need input scaling, early stopping and
      several seeds; Random Forest is simpler to train and explain. For a
      quick, robust tool, Random Forest remains a strong choice; for best
      accuracy, a small neural network.
  \end{itemize}
\end{frame}

\section{Real Weather: NASA POWER}

\begin{frame}{Replacing synthetic weather with real data}
  \small
  \begin{itemize}
    \item \textbf{NASA POWER}: daily satellite-and-reanalysis weather from NASA's
      Agroclimatology service, free, no account needed.
    \item Same location (21.15$^\circ$N, 79.09$^\circ$E, Nagpur), same 2,000
      days (1 Jan 2019 -- 22 Jun 2024), \textbf{0 missing values}.
  \end{itemize}
  \begin{table}
    \centering
    \footnotesize
    \begin{tabular}{ll}
      \toprule
      \textbf{Model input} & \textbf{NASA POWER variable} \\
      \midrule
      Rainfall & PRECTOTCORR (bias-corrected, based on NASA GPM satellite rain) \\
      Temperature & T2M (plus real daily max / min for the FAO-56 calculation) \\
      Humidity & RH2M \\
      Wind speed & WS2M (at 2 m height, as FAO-56 requires) \\
      Solar radiation & ALLSKY\_SFC\_SW\_DWN \\
      \bottomrule
    \end{tabular}
  \end{table}
  \begin{itemize}
    \item The irrigation demand is still \textbf{calculated} with the same FAO-56
      method, maize crop calendar and settings --- only the weather is now real.
  \end{itemize}
\end{frame}

\begin{frame}{Real weather vs.\ the synthetic weather}
  \centering
  \fitfig[0.6]{figures/13_nasa_vs_synthetic_weather.png}

  \footnotesize
  The synthetic data had the seasons right. Real weather differs in the details:
  monsoon \textbf{wind is much calmer} (about 2.6 vs 4.5 m/s), rain is
  \textbf{more persistent} from day to day (correlation 0.63 vs 0.27), and
  pre-monsoon days are hotter and drier, so on 199 days atmospheric thirst
  exceeds the 8 mm/day cap used in the synthetic data (peak 13.8 mm/day).
  That cap was therefore removed.
\end{frame}

\begin{frame}{Classical models on real weather}
  \centering
  \fitfig[0.56]{figures/14_nasa_classical_models.png}

  \small
  \begin{itemize}
    \item Same ranking as on synthetic weather: Random Forest best
      (0.42 mm/day, R\textsuperscript{2} 0.971), Linear Regression worst
      (0.91 mm/day, 13 negative predictions).
    \item Test period: the last 400 days, 20 May 2023 -- 22 Jun 2024.
  \end{itemize}
\end{frame}

\begin{frame}{Classical ML vs.\ deep learning on real weather}
  \begin{table}
    \centering
    \small
    \begin{tabular}{llccc}
      \toprule
      & \textbf{Model} & \textbf{MAE} & \textbf{RMSE} & \textbf{R\textsuperscript{2}} \\
      \midrule
      \multirow{4}{*}{\rotatebox{90}{\footnotesize Classical}}
      & Linear Regression & 0.91 & 1.30 & 0.880 \\
      & Decision Tree & 0.69 & 0.97 & 0.933 \\
      & Random Forest & 0.42 & 0.64 & 0.971 \\
      & Random Forest + 7-day window & 0.49 & 0.77 & 0.959 \\
      \midrule
      \multirow{3}{*}{\rotatebox{90}{\footnotesize Deep}}
      & MLP (today only) & 0.27 {\scriptsize$\pm$0.02} & 0.36 & 0.991 \\
      & 1D-CNN (7-day) & 0.33 {\scriptsize$\pm$0.02} & 0.46 & 0.985 \\
      & \textbf{LSTM (7-day)} & \textbf{0.23} {\scriptsize$\pm$0.01} & \textbf{0.29} & \textbf{0.994} \\
      \bottomrule
    \end{tabular}
  \end{table}
  \footnotesize Same 7 weather and crop inputs; deep models are the mean of 3 seeds.
  \small
  \begin{itemize}
    \item Same ranking as on synthetic data. The LSTM cuts Random Forest's error by
      \alert{47\%} (0.42 $\rightarrow$ 0.23 mm/day) and does slightly
      \emph{better} on real weather than on synthetic (0.26).
    \item Real rain and heat spells last several days, so the 7-day memory now
      helps more: LSTM beats the today-only MLP by 0.04 mm/day (was 0.01).
  \end{itemize}
\end{frame}

\begin{frame}{All models at a glance, real weather}
  \centering
  \fitfig[0.72]{figures/15_nasa_ml_vs_dl.png}
\end{frame}

\begin{frame}{Predicted vs.\ actual, real weather}
  \centering
  \fitfig[0.6]{figures/16_nasa_actual_vs_predicted.png}

  \small Best tree model vs.\ best network on the last 400 real days.
\end{frame}

\begin{frame}{What the model relies on, real weather}
  \centering
  \fitfig[0.56]{figures/17_nasa_feature_importance.png}

  \small
  \begin{itemize}
    \item Crop coefficient still leads. \textbf{Temperature} rises to second:
      real hot, dry pre-monsoon days drive the highest demand.
    \item Wind and humidity matter more than in the synthetic data
      (permutation importance about 2--4$\times$ higher).
  \end{itemize}
\end{frame}

\begin{frame}{Overfitting check, real weather}
  \begin{table}
    \centering
    \small
    \begin{tabular}{ccccc}
      \toprule
      \textbf{LSTM seed} & \textbf{Epochs} & \textbf{Training days} & \textbf{Validation days} & \textbf{Test days} \\
       & & \footnotesize MAE / R\textsuperscript{2} & \footnotesize MAE / R\textsuperscript{2} & \footnotesize MAE / R\textsuperscript{2} \\
      \midrule
      0 & 179 & 0.16 / 0.996 & 0.22 / 0.993 & 0.23 / 0.994 \\
      1 & 233 & 0.18 / 0.996 & 0.20 / 0.994 & 0.23 / 0.994 \\
      2 & 139 & 0.17 / 0.996 & 0.21 / 0.994 & 0.22 / 0.994 \\
      \bottomrule
    \end{tabular}
  \end{table}
  \vspace{-0.6em}
  \footnotesize
  \begin{itemize}
    \item \textbf{LSTM: no meaningful overfitting.} Test error is only
      $\approx$0.06 mm/day above training error, and test R\textsuperscript{2}
      (0.994) is as high as training.
    \item Training error stays above the $\approx$0.14 mm/day noise floor in the data.
    \item The \textbf{1D-CNN} is the one model with a visible gap (training 0.16
      vs test 0.33 mm/day): mild overfitting, which is why it ranks last of the
      three networks.
  \end{itemize}
\end{frame}

\begin{frame}{Synthetic vs.\ real weather: summary}
  \begin{table}
    \centering
    \small
    \begin{tabular}{lcc}
      \toprule
      \textbf{7 weather and crop inputs} & \textbf{Synthetic} & \textbf{NASA POWER} \\
      \midrule
      Random Forest, MAE / R\textsuperscript{2} & 0.40 / 0.963 & 0.42 / 0.971 \\
      MLP, MAE / R\textsuperscript{2} & 0.28 / 0.990 & 0.27 / 0.991 \\
      LSTM, MAE / R\textsuperscript{2} & \textbf{0.26 / 0.990} & \textbf{0.23 / 0.994} \\
      \bottomrule
    \end{tabular}
  \end{table}
  \begin{itemize}
    \item The conclusions \alert{hold on real weather}: the weather and crop
      inputs alone are enough, and the LSTM is the most accurate model.
  \end{itemize}
\end{frame}

\section{Wrapping Up}

\begin{frame}{Limitations}
  \begin{itemize}
    \item The main results use synthetic weather. The real-weather test uses
      NASA POWER, but the irrigation demand is still \emph{calculated} with
      FAO-56, not measured in a real, instrumented field.
    \item NASA POWER is gridded satellite/reanalysis data (cells tens of km
      wide), not a weather station inside the field.
    \item Only one crop (maize) and one climate region were modelled.
    \item A simplified rainfall formula was applied at a daily step, though it
      was originally meant for monthly totals.
    \item The models predict today's demand from today's (and the past
      week's) weather; multi-day-ahead forecasting would need weather
      \emph{forecasts} as inputs.
    \item The target follows the FAO-56 formula exactly, which favours smooth
      models such as neural networks; measured irrigation would be noisier.
  \end{itemize}
\end{frame}

\begin{frame}{Conclusion}
  \begin{itemize}
    \item Machine learning can predict daily irrigation water demand very
      accurately from simple weather and crop inputs alone: Random Forest
      reaches R\textsuperscript{2} = 0.96.
    \item Deep learning goes further: an \textbf{LSTM} reaches
      R\textsuperscript{2} = 0.99 and cuts the typical error by 35\% (0.26
      mm/day), and a simple MLP is nearly as good.
    \item The results \textbf{hold on real NASA POWER weather}: LSTM
      R\textsuperscript{2} = 0.994 (0.23 mm/day), with no meaningful overfitting.
    \item Results consistently make physical sense: right seasonal pattern,
      right feature priorities (cross-checked several ways), no
      systematic bias.
    \item This points toward a faster, simpler decision-support tool for
      irrigation scheduling, without sacrificing physical grounding.
  \end{itemize}
\end{frame}

{
\setbeamertemplate{footline}{}
\begin{frame}[standout]
  Thank You \\
  \vspace{0.5em}
  {\normalsize Questions?}
\end{frame}
}

\end{document}